\chapter*{Conclusion générale}
\addcontentsline{toc}{chapter}{Conclusion générale}
\markboth{\textbf{Conclusion générale}}{}

Le travail réalisé dans le cadre de ce projet de fin d'études a abouti au développement d'AFYA, une application mobile dédiée à la gestion intelligente des médicaments et au suivi des traitements. Cette solution répond à un besoin réel et quotidien des familles et des équipes soignantes, en facilitant l'organisation de la pharmacie familiale, en améliorant le suivi thérapeutique et en renforçant la coordination entre les membres d'un même groupe. L'application permet ainsi de réduire les risques liés à l'oubli de prise de médicaments, à la consommation de produits périmés et au manque de visibilité sur les stocks disponibles.

Le développement a été organisé de manière progressive et structurée autour de trois versions majeures. La première version a établi les fondations essentielles de l'application en mettant en place les mécanismes de création de compte, d'accès sécurisé et de gestion des groupes familiaux ou médicaux. Cette étape a également introduit les premières fonctionnalités permettant d'ajouter des médicaments à sa pharmacie personnelle et de consulter les informations associées. La deuxième version a enrichi l'application avec des fonctionnalités avancées de gestion des traitements, d'enregistrement des prescriptions et de configuration de rappels de prise de médicaments. Ces mécanismes visent à améliorer l'observance thérapeutique grâce à des notifications opportunes et un suivi détaillé de l'historique des prises. Enfin, la troisième version a renforcé la dimension collaborative de l'application en permettant aux membres d'un groupe de partager les informations médicales, de gérer collectivement la pharmacie et d'assurer un suivi coordonné des traitements de leurs proches. Chaque version a été livrée de manière incrémentale, permettant ainsi de valider progressivement les fonctionnalités développées et d'intégrer les retours d'amélioration.

Ce projet a constitué une expérience formatrice riche en apprentissages, tant sur le plan de la conception d'applications mobiles que sur celui de la gestion de projet. Il nous a permis de nous confronter à des problématiques réelles du domaine de la santé et du bien-être, d'adopter une méthodologie de travail structurée et itérative, et de développer une solution concrète répondant aux attentes des utilisateurs finaux. Le respect des normes de sécurité, la prise en compte de l'expérience utilisateur et la coordination entre les différentes fonctionnalités ont été au cœur de nos préoccupations tout au long du développement.

En perspective, l'application AFYA offre de nombreuses possibilités d'enrichissement futur. Parmi les évolutions envisageables figurent l'intégration de fonctionnalités d'analyse des interactions médicamenteuses, l'ajout de recommandations personnalisées basées sur l'historique de santé, le développement d'un système de téléconsultation ou encore l'extension vers d'autres plateformes pour toucher un public plus large. Ces améliorations permettraient de renforcer encore davantage l'utilité de l'application et son impact positif sur la santé et le bien-être des utilisateurs.
